% LSPV-1 REPORT::MAJORANTNOGO1::2026-10-12T09:44:39Z
% !TEX program = xelatex
\documentclass[11pt]{article}
\usepackage[margin=0.86in]{geometry}
\usepackage{amsmath,amssymb,amsthm,booktabs,array,enumitem,hyperref}
\hypersetup{colorlinks=true,linkcolor=blue,urlcolor=blue}
\newtheorem{theorem}{Theorem}
\newtheorem{corollary}{Corollary}
\newcommand{\OPEN}{\textbf{OPEN}}
\newcommand{\EXACT}{\textbf{PROVED/EXACT}}
\newcommand{\KNOWN}{\textbf{KNOWN THEOREM}}
\newcommand{\COMP}{\textbf{COMPUTATIONAL}}
\newcommand{\RETIRED}{\textbf{FALSIFIED/RETIRED}}
\newcommand{\same}{\textsc{same\_realization}}
\title{MAJORANTNOGO1: No Polynomial Faithful Resource Majorant\\Survives Equivalent Polynomial Padding}
\author{LSPV-1 / UNIVERSAL\_CORE\_LOWER\_BOUND}
\date{2026-10-12 09:44 UTC}

\begin{document}
\maketitle

\begin{center}
\fbox{\parbox{0.94\linewidth}{
\textbf{Executive status.} P versus NP remains \OPEN.  This run tests the
obvious repair left after exact time and exact joint time--space were shown to
admit padding gauges: replace presentation-dependent costs by a resource
majorant that is invariant under every equivalent polynomial presentation.
That repair is impossible if the majorant must also dominate the actual source
time or space and have one fixed polynomial bound.  The impossibility theorem
is \EXACT; a 787,320-assertion finite hostile-control audit passed with zero
failures and is \COMP.  The result is a no-go theorem, not a SAT lower bound.
}}
\end{center}

\section{Claim card and exact target}

\paragraph{Hard clean core.}
Let $\mathcal H^{\pm}_{2,2}$ be the duplicate-free proper sign-monochromatic
$3$-CNFs in which every variable occurs exactly twice positively and twice
negatively.  If a formula has $n$ variables and $m$ clauses, then $3m=4n$.
The equality is \EXACT, and NP-completeness of satisfiability on this family is
a \KNOWN{} theorem of D\"ocker.  The encoded input length is denoted $N$.

\paragraph{Machine and charges.}
Fix a standard deterministic multitape machine model.  For a total exact
one-bit decider $A$, let
\[
 \tau_A(N)=\max_{|x|=N}T_A(x),\qquad
 \sigma_A(N)=\max_{|x|=N}S_A(x),
\]
where input access, simulation, counters, wrapper recognition, dummy
transitions, allocation, touching and retention of work cells, and output are
all charged.  Extraction, reconstruction and communication are absent from
the wrapper below except for storing and returning the one decision bit.

\paragraph{Candidate under test.}
A tempting repair is to postulate resource functionals $U_T,U_S$ on program
presentations satisfying simultaneously:
\begin{enumerate}[label=(\roman*),leftmargin=1.8em]
\item \emph{padding invariance}: $U_R(A,N)=U_R(J_{q,r}(A),N)$ for every fixed
      $q,r$ and $R\in\{T,S\}$;
\item \emph{faithfulness}: $U_T(A,N)\ge\tau_A(N)$ and
      $U_S(A,N)\ge\sigma_A(N)$; and
\item \emph{one polynomial envelope}: there are fixed constants
      $C,D,c,d,N_0$ such that $U_T(A,N)\le C N^c$ and
      $U_S(A,N)\le D N^d$ for every polynomial-time presentation $A$ and
      every $N\ge N_0$.
\end{enumerate}
Full replacement invariance is stronger than (i), so falsifying (i)--(iii)
also falsifies the full repair.

\section{Primary-literature and barrier-transfer preflight}

The two blocking obstacles are (a) presentation dependence under arbitrary
fixed polynomial padding and (b) model-matched landing: a structural or
semantic quantity does not itself lower-bound an unrestricted decider.
Primary literature was checked before the new claim.  Hartmanis--Stearns
establish the model-sensitive hierarchy viewpoint; Blum gives axioms for
machine-independent complexity measures but does not provide the invariant
majorant proposed here; Fortnow--Lipton--van Melkebeek--Viglas (FLMV) gives the
relevant unconditional random-access SAT time--space endpoint.  The present
wrapper theorem is elementary and does not invoke Blum speedup.

Baker--Gill--Solovay and Aaronson--Wigderson remain active audits.  The padding
construction relativizes and algebrizes, so it diagnoses a failed route and
cannot be advertised as the nonrelativizing hinge needed for $P\ne NP$.
Razborov--Rudich is not directly triggered because no constructive large
circuit property is asserted.  The exact transfer classification is:
\begin{center}
\begin{tabular}{p{0.28\linewidth}p{0.18\linewidth}p{0.45\linewidth}}
\toprule
Source & Classification & Reason \\
\midrule
Blum complexity axioms & useful structural analogy & Program-indexed measures
remain presentation sensitive; no theorem yields (i)--(iii). \\
FLMV SAT endpoint & directly transferable theorem & Only after a single
ordinary machine is proved to have absolute model-matched exponents. \\
VLSI area--time, routing, and memory hierarchy & misleading without a bridge &
Those lower bounds constrain declared layouts/models, not arbitrary SAT
program replacements. \\
Control observability / barrier certificates & analogy/candidate & They require
a declared state system; arbitrary one-bit SAT computation need not expose it. \\
\bottomrule
\end{tabular}
\end{center}

\section{Exact symbolic falsification}

For fixed integers $q,r\ge1$, define $J_{q,r}(A)$ as follows on input $x$ of
length $N$: run $A(x)$ to completion, store its bit $b$, execute at least
$N^q$ real dummy transitions, touch and retain at least $N^r$ work cells
(charging the touching time), and return $b$.  The wrapper syntax contains
$q,r$ and the embedded code for $A$, so a fixed parser strips it without
solving SAT.

\begin{theorem}[\texttt{SEMANTIC-RESOURCE-MAJORANT-NOGO1}]
Let $L$ be any language having at least one deterministic polynomial-time
decider on the fixed model, and let $\mathsf{PolyDec}(L)$ be all such
presentations.  There do not exist functionals $U_T,U_S$, constants
$C,D>0$, integers $c,d$, and $N_0$ satisfying (i)--(iii) above.  Indeed, the
conclusion already follows if padding invariance is required only for the
family $J_{q,r}$.
\end{theorem}

\begin{proof}
Fix $A\in\mathsf{PolyDec}(L)$.  For every fixed $q,r$, the machine
$J_{q,r}(A)$ decides exactly $L$ and remains polynomial-time; touching
$N^r$ cells contributes polynomial time and peak workspace.  Moreover,
\[
 \tau_{J_{q,r}(A)}(N)\ge N^q,
 \qquad
 \sigma_{J_{q,r}(A)}(N)\ge N^r.
\]
Choose $q=c+1$ and $r=d+1$.  Padding invariance and faithfulness imply
\[
 U_T(A,N)=U_T(J_{q,r}(A),N)\ge N^{c+1},\qquad
 U_S(A,N)=U_S(J_{q,r}(A),N)\ge N^{d+1}.
\]
For every $N>\max\{C,D,N_0\}$ these inequalities contradict respectively
$U_T(A,N)\le C N^c$ and $U_S(A,N)\le D N^d$.  No semantic property of $L$
was used.  Thus the theorem applies in particular to SAT and to any
hypothetical polynomial decider for $\mathcal H^{\pm}_{2,2}$.
\end{proof}

\begin{corollary}[exact quantifier form]
The following prefix is false for every language $L\in P$:
\[
\begin{split}
\exists U_T,U_S,C,D,c,d,N_0\ \forall A\in\mathsf{PolyDec}(L)\
\forall N\ge N_0:\quad
&\tau_A(N)\le U_T(A,N)\le C N^c,\\
&\sigma_A(N)\le U_S(A,N)\le D N^d,
\end{split}
\]
when $U_T,U_S$ are invariant under all $J_{q,r}$ presentations.  Therefore
\texttt{WRAPPER-INVARIANT-FAITHFUL-POLY-MAJORANT1} is \RETIRED.
\end{corollary}

The theorem does \emph{not} rule out a wrapper-invariant lower bound, an
output-specific absolute resource theorem, or an invariant that is not a
majorant of raw source resources.  It rules out only the conjunction actually
quantified above.

\section{Bounded reproducible hostile-control audit}

The deterministic script
\path{experiments/review/semantic_resource_majorant_nogo_audit.py}
wrote
\path{results/raw/review/2026-10-12_semantic_resource_majorant_nogo.json}.
One invocation completed in $0.145$ seconds, so no chunking was needed.  It
checked exponents $c,d=1,\ldots,9$, constants
$1,2,3,5,8,13$, nine families, five encodings and six presentations.  There
were 131,220 rows and 787,320 exact integer assertions, with zero failures.

The families were signed clean core, Horn, 2SAT, affine/Tseitin,
bounded-treewidth, independent products, padded formulas, redundant-clause
formulas, and a unit contradiction.  Encodings were canonical, renamed,
complemented, auxiliary-variable and padded; presentations were
DIRECT/COPY/DEAD/INTERP/RECOMP/PREPAD.  Preprocessing was included in the
baseline, wrapper touching was included in time, retained cells were included
in peak workspace, and no extraction or reconstruction was performed.  This
finite audit is \COMP{} only; the theorem is supported by the proof above, not
by extrapolating the audit.

\section{Semantic complexity versus actual algorithmic cost}

The construction leaves the input, decision language, witness space,
compatibility relation $H(F,\alpha)$, branching-program size, circuit
representation and one-bit output unchanged.  It changes only real executed
work.  Conversely, Horn propagation, implication-graph SCCs, Gaussian
elimination, bounded-treewidth dynamic programming and compact alternate
representations can solve their promised subclasses without materializing a
chosen global compatibility object.  Thus neither large semantic-state count
nor low semantic complexity controls the wrapper theorem.  This is exactly why
a source-resource majorant cannot serve as an intrinsic language invariant.

\section{SAME\_REALIZATION and smallest repair}

The parser $J_{q,r}(A)\mapsto A$ is an external syntax operation.  It does not
show that the padded execution internally exposed, reconstructed or paid for a
clean-core object.  Treating the stripped program's resources as those of the
executed padded program would violate \same.  Treating an invariant majorant
as if it dominated every equivalent source presentation is now ruled out by
the theorem.

The smallest surviving repair is output-specific:
\texttt{SELECTED-OUTPUT-ABSOLUTE-LANDING1} (\OPEN).  A charged uniform
transformation may select one realized output machine $B=C(A)$ and prove
absolute encoded-length bounds for that very $B$ in the cited machine model.
It must charge construction, parsing, source access, calls, replay, OR,
workspace and output, and it must not claim that the bound majorizes all
programs equivalent to $B$.  A separate \same{} lemma must identify this
realization with the FLMV endpoint.  CALQNT1 and TDLG1 show that universally
landing every hypothetical polynomial SAT decider in a forbidden region is
already separation-strength; the present theorem supplies no such landing.

\section{Post-quantum technique ledger}

Schur squares, quadratic-relation ranks, Hilbert functions and syzygies are
\textbf{ANALOGY/CANDIDATE needing a new bridge}: they may distinguish a
chosen representation but do not survive arbitrary program replacement as
faithful source-cost majorants.  Collision entropy and
$|\operatorname{im}L|\ge |A|^2/\operatorname{Coll}(L)$ are exact counting
tools at a declared proof-to-node map, but there is no such map in this
one-bit wrapper; transfer is \textbf{FAILED TRANSFER}.  Worst-case/average-case
extraction and special-soundness analogies remain candidates only after a
black-box extraction theorem; hashes and random oracles are inadmissible for
an unconditional separation.  In every case, exploitable hidden algebraic
structure may aid an algorithm rather than certify hardness.

\section{CORPUS REVIEWED}

The current GitHub tree was surveyed through repository search.  Direct reads
included root \path{main.tex}, \path{AGENTS.md}, \path{CODEX.md},
\path{.codex/skills/pnp-research/SKILL.md}, actual \path{main4.tex} and
\path{main5.tex}, all seven task specifications, newest verified
MAIN/SYNTHESIS-0023, TDLG1, KCALIB1, COVCOLL1, CLKSLACK1, TIGHTDELAY1,
CALQNT1, JOINTGAUGE1 and MODELBRIDGE-0010.  Relevant historical HEC,
VSC-HEC, G-PPC-HEC, CRS-PPC and CCSPC causal components, governance, and
cited experiment/result paths were inspected rather than inferred from peer
summaries.  The newest verified predecessor LSPV-5 write was
\path{research/main/2026-10-12_0850_MAIN_SYNTHESIS-0023.tex}; its header and
sentinel were present, so no recovery retry was needed.

\section{DUPLICATION / CONTRADICTION AUDIT}

TIGHTDELAY1 proves exact-time padding; JOINTGAUGE1 proves simultaneous
time/space padding; CLKSLACK1 concerns loose declared clocks;
MODELBRIDGE-0010 rules out a fixed-exponent normalizer by a hierarchy
argument; COVCOLL1 concerns decision-fiber collapse; TDLG1 and CALQNT1 expose
the forbidden-endpoint quantifiers.  None states the present trilemma for a
replacement-invariant \emph{faithful majorant}.  The new theorem is an exact
corollary of the padding construction, not a new padding construction and not
Blum speedup.  It contradicts only the proposed majorant repair, not any
verified corpus theorem.

\section{WHAT CHANGED RELATIVE TO ROOT MAIN}

Root MAIN already records that relative exact time and joint time/space
contractions are presentation gauges.  This run adds the stronger boundary:
the suggested absolute repair cannot be defined as a polynomially bounded,
wrapper-invariant upper envelope that still faithfully dominates every
equivalent source program.  \texttt{ABSOLUTE-JOINT-CROSSREGIME1} must therefore
refer to one selected output realization or to a non-majorant invariant with a
separate transfer theorem.

\section{ADOPTED / REJECTED / RETIRED PRIOR COMPONENTS}

\begin{description}[leftmargin=2.5em]
\item[ADOPTED.] The exact $J_{q,r}$ padding wrapper from JOINTGAUGE1, the
  signed hard-core definition, TDLG1's total-domain guard, CALQNT1's quantifier
  lattice, and the model-matched FLMV endpoint.
\item[REJECTED.] Semantic-state counting, proof-system/representation lower
  bounds without arbitrary-decider transfer, and restricted online/VLSI/
  communication analogies as unrestricted SAT lower bounds.
\item[RETIRED.] \texttt{WRAPPER-INVARIANT-FAITHFUL-POLY-MAJORANT1}.
\item[OPEN.] A charged selected-output construction with absolute time--space
  coordinates, \same{} transfer, and a nonrelativizing/nonalgebrizing hinge.
\end{description}

\section{STALE OR OVERSTATED CLAIMS FOUND}

The phrase ``absolute, wrapper-invariant law whose output time and space are
bounded by encoded input length'' is safe only if ``output'' means the one
selected realized machine.  If read as a faithful polynomial upper majorant
for every equivalent presentation, it is false by this report.  Root MAIN did
not explicitly assert that stronger reading, so the update narrows wording
rather than retracting a verified theorem.

\section{Strict status ledger and handoff}

\begin{center}
\begin{tabular}{p{0.55\linewidth}p{0.26\linewidth}}
\toprule
Item & Status \\
\midrule
\texttt{SEMANTIC-RESOURCE-MAJORANT-NOGO1} & \EXACT \\
Finite hostile-control inequality audit & \COMP \\
FLMV model-matched SAT time--space endpoint & \KNOWN \\
\texttt{WRAPPER-INVARIANT-FAITHFUL-POLY-MAJORANT1} & \RETIRED \\
\texttt{SELECTED-OUTPUT-ABSOLUTE-LANDING1} & \OPEN \\
Arbitrary-decider clean-core transfer & \OPEN \\
P versus NP & \OPEN \\
\bottomrule
\end{tabular}
\end{center}

\paragraph{Precise missing universal implication.}
Given an arbitrary uniform polynomial-time SAT decider $A$, construct and
fully charge one specific ordinary machine $B=C(A)$ whose absolute encoded-
length time and peak-space exponents lie in a cited unconditionally forbidden
region, then prove that those bounds apply to that same realized $B$.  The
majorant theorem shows that this implication cannot be discharged by assigning
one faithful polynomial resource envelope to every equivalent presentation.

\begin{thebibliography}{9}
\bibitem{HS65}
J. Hartmanis and R. E. Stearns,
``On the Computational Complexity of Algorithms,''
\emph{Transactions of the AMS} 117 (1965), 285--306.
\url{https://doi.org/10.1090/S0002-9947-1965-0170805-7}.

\bibitem{Blum67}
M. Blum, ``A Machine-Independent Theory of the Complexity of Recursive
Functions,'' \emph{JACM} 14(2) (1967), 322--336.
\url{https://doi.org/10.1145/321386.321395}.

\bibitem{FLMV05}
L. Fortnow, R. Lipton, D. van Melkebeek and A. Viglas,
``Time--Space Lower Bounds for Satisfiability,'' \emph{JACM} 52(6) (2005).
\url{https://pages.cs.wisc.edu/~dieter/Papers/sat-ts-jacm.ps}.

\bibitem{BGS75}
T. Baker, J. Gill and R. Solovay, ``Relativizations of the $P=?NP$
Question,'' \emph{SIAM Journal on Computing} 4(4) (1975).
\url{https://doi.org/10.1137/0204037}.

\bibitem{RR97}
A. Razborov and S. Rudich, ``Natural Proofs,'' \emph{Journal of Computer and
System Sciences} 55(1) (1997), 24--35.
\url{https://doi.org/10.1006/jcss.1997.1494}.

\bibitem{AW08}
S. Aaronson and A. Wigderson, ``Algebrization: A New Barrier in Complexity
Theory,'' ECCC TR08-005 (2008).
\url{https://eccc.weizmann.ac.il/report/2008/005/}.
\end{thebibliography}

\noindent\textbf{Final status:} P versus NP remains \OPEN.

% END_OF_VERIFIED_REPORT::LSPV-1::MAJORANTNOGO1
\end{document}
